\documentclass[10pt,a4paper]{amsart}
\usepackage[margin=19mm]{geometry}
\usepackage{amsmath,amssymb,mathtools}
\usepackage[scr=rsfs,cal=euler]{mathalfa}
\usepackage{xfrac}
\usepackage[unicode,hidelinks,bookmarks=false]{hyperref}
\newcommand{\ie}{i.e.\ }
\newcommand{\cf}{cf.\ }
\newcommand{\resp}{respectively\ }
\newcommand{\Subsection}{\subsection}
\providecommand{\nobreakdash}{\nobreak\hbox{-}}
\setlength{\emergencystretch}{2em}
\multlinegap0pt
\usepackage{bm}
\usepackage{bbm}
\usepackage{dsfont}
\usepackage{xspace}
\usepackage{cancel}
\usepackage{mathtools}
\usepackage{amsthm}
\makeatletter
\@ifundefined{theorem}{\newtheorem{theorem}{Theorem}}{}
\@ifundefined{lemma}{\newtheorem{lemma}[theorem]{Lemma}}{}
\makeatother
\newcommand{\lk}[2]{\operatorname{lk}(#1,#2)} %link
\newcommand{\pcd}{\operatorname{pcd}}
\newcommand{\st}[2]{\operatorname{st}(#1,#2)} %star
\renewcommand{\geq}{\geqslant}
\renewcommand{\setminus}{\smallsetminus}
\title[Convex cocompact groups in real hyperbolic spaces with limit]{Convex cocompact groups in real hyperbolic spaces with limit set a Pontryagin sphere}
\author{Sami Douba, Gye-Seon Lee, Ludovic Marquis, Lorenzo Ruffoni}
\date{Source: arXiv:2502.09470v3 (CC BY 4.0)}
\begin{document}
\maketitle

\noindent\emph{Source and licence.} Convex cocompact groups in real hyperbolic spaces with limit set a Pontryagin sphere. Version arXiv:2502.09470v3 (CC BY 4.0), distributed under the Creative Commons Attribution 4.0 International licence, as recorded in the arXiv licence metadata for this exact version and shown on the abstract page https://arxiv.org/abs/2502.09470v3. Full rights notice: licenses/arxiv-2502.09470-CC-BY-4.0.txt.\par\smallskip

\noindent\textit{Editorial scope.} Counted unit: the Lemma whose source label is \texttt{lem:menger\_vertex} in the article (source0.tex lines 1152--1158, source label \texttt{lem:menger\_vertex}), delivered together with the article's own complete original proof of it (source0.tex lines 1159--1182, one \texttt{proof} environment of 24 lines). No mathematics is written, repaired or re-derived: statement and proof are the article's own text line for line. Required context retained uncounted: none. Every definition, display and label of those lines is retained unchanged. Omitted from this package and not redistributed: 1--1151, 1183--1272. Nothing inside a retained excerpt is omitted or altered: every retained body is delivered exactly as the article's lines are written, so no omission receipt of any kind accompanies this package. Citations: the retained excerpts carry no citation command at all, so no bibliography entry of the article is transcribed and no reference is redistributed. Reference adaptation: none was needed and none was made; every reference inside the delivered unit resolves inside it, so no unresolved reference or citation is produced. Printed slips kept literal: no printed word, symbol, hypothesis or number of the counted statement or of its proof was altered. Layout and class: the article's own class is \texttt{amsart} with packages including inputenc, babel, csquotes, color, hyperref, biblatex, graphicx, euscript, mathrsfs, tikz, todonotes, fouriernc, helvet, courier; the wrapper is the lane's pinned \texttt{amsart} editorial wrapper at 10pt on A4, which restates the theorem-like environments the retained text needs and adds the article's own macro definitions that the retained text needs (\texttt{\textbackslash geq}, \texttt{\textbackslash lk}, \texttt{\textbackslash pcd}, \texttt{\textbackslash setminus}, \texttt{\textbackslash st}), with extra wrapper packages (bm, bbm, dsfont, xspace, cancel, mathtools). The delivered numbering is the unit's own. Assets: no retained span contains a figure environment, a graphics-inclusion command, a PDF-page inclusion, a table or a TikZ picture; no image file is copied into this package, the package declares no asset dependency and no image file is redistributed with it. Licence: version arXiv:2502.09470v3 of this article is distributed under the Creative Commons Attribution 4.0 International licence, as recorded in the arXiv licence metadata for this exact version and shown on the abstract page https://arxiv.org/abs/2502.09470v3. A full rights notice is delivered at licenses/arxiv-2502.09470-CC-BY-4.0.txt. Characters: every retained excerpt is pure ASCII, so the delivered engine, XeLaTeX with its default Latin Modern text font, reports no missing character.
\begin{lemma}\label{lem:menger_vertex}
    Let ${K}$ be a flag-no-square triangulation  of a  closed connected orientable surface $S$ of genus $g\geq 1$.
    Let $M \subseteq {K}$ be a non-planar inseparable full subcomplex. 
    Let $v \in M^{(0)}$ such that $\lk{v}{M}$ is a circle,
    and let $L$ be the full subcomplex of $M$ spanned by $M^{(0)}\setminus \{v\}$.
    Then $L$ is a non-planar inseparable full subcomplex with $\pcd (L)=1$.
\end{lemma}

\begin{proof}
We start by noticing that both $M$ and $L$ are   flag-no-squares   complexes of dimension at most $2$, because they are full subcomplexes of $ K$. 
Since $M$ is non-planar and not a sphere, and $L$ is obtained from $M$ by removing the open disk bounded by the circle $\lk vM$, it follows that $L$ is non-planar as well.
 
To compute the cohomological dimension, note that for any (possibly empty) simplex $\sigma$ of $L$, the complex $L\setminus \sigma$ deformation retracts to a graph (it can be thickened to a subsurface with non-empty boundary of $S$).
In particular, the cohomology of $L\setminus \sigma$ always vanishes in degree $2$ and above.
On the other hand, note that $L$ does not deformation retract to a tree, because $L$ is non-planar. So we conclude that $\pcd(L)=1$.

To conclude, we show that $L$ is inseparable. 
Since $M$ is connected and $\lk vM$ is connected, we have that $L$ is connected. 
(This follows from Mayer--Vietoris applied to the decomposition  $M=\st vM \cup L$, for which $\st vM\cap L=\lk vM$.)  
Let $A\subseteq L$ be a subcomplex that disconnects $L$ but not $M$.
We need to check that~$A$ is not among the forbidden subcomplexes.
Since $\lk vM$ is connected, $A\cap \lk vM$ must disconnect $\lk vM$.
Since $\lk vM$ is a circle, $A\cap \lk vM$ must have at least $2$ non-adjacent vertices.
In particular, $A$ cannot be a simplex.
For the remaining cases we argue by contradiction.
Let $A=\{p,q\}\ast \sigma$ be the suspension of a simplex. Then $A\cap \lk vM$  consists of the two suspension points $\{p,q\}$. 
But then for any vertex $r\in \sigma$ we have that $p,q,r,v$ form an induced square in $M$, which cannot happen.
 Finally, let $A=\{u,w\}\subseteq \lk vM$ be a pair of non-adjacent vertices. 
 But then $\{u,v,w\}$ separates $M$.
 This is absurd since $M$ has no separating suspension of a simplex.
 This concludes the proof that $L$ is inseparable. 
\end{proof}

\end{document}
